% Part: counterfactuals
% Chapter: minimal-change-semantics
% Section: transitivity

\documentclass[../../../include/open-logic-section]{subfiles}

\begin{document}

\olfileid{cnt}{min}{tra}

\olsection{পরিযায়িতা}

বস্তুগত শর্তবচনের জন্য শৃঙ্খল-নিয়মটি চলে:
$!A \lif !B, !B \lif !C \Entails !A \lif !C$। অন্য কথায়,
বস্তুগত শর্তবচন পরিযায়ী। প্রতিবাস্তব শর্তবচনের ক্ষেত্রেও
কি তা সত্য? স্টালনেকারের দেওয়া নিচের উদাহরণটি দেখুন।
\begin{quote}
  জে.~এডগার হুভার যদি রাশিয়ায় জন্মাতেন, তবে কমিউনিস্ট হতেন।

  জে.~এডগার হুভার যদি কমিউনিস্ট হতেন, তবে বিশ্বাসঘাতক হতেন।

  অতএব, জে.~এডগার হুভার যদি রাশিয়ায় জন্মাতেন, তবে বিশ্বাসঘাতক হতেন।
\end{quote}
% BN-SRC-429: হুভারের জন্ম ১৮৯৫; সোভিয়েত ইউনিয়নের প্রতিষ্ঠা ১৯২২, তাই জন্ম রাশিয়ায়।
হুভার যদি যুক্তরাষ্ট্রের বদলে রাশিয়ায় জন্মাতেন (বাস্তবে যে সময়ে
জন্মেছিলেন, সেই সময়েই), তবে রাশিয়ায় বড় হয়ে পরে সোভিয়েত
ইউনিয়নে কমিউনিস্ট হতেন—এমনটা ধরা যাক। তাই প্রথম পূর্বধারণাটি সত্য। দ্বিতীয়
পূর্বধারণাটিও উদাহরণটির অনুমান মেনে আলাদা করে বিবেচনা করলে
সত্য ধরা হচ্ছে। কিন্তু সিদ্ধান্তটি
মিথ্যা:
% BN-SRC-429: দ্বিতীয় জন্মস্থান-উল্লেখেও ১৮৯৫ সালের রাশিয়া; সোভিয়েত ইউনিয়ন তখন গঠিত হয়নি।
হুভার রাশিয়ায় জন্মালে পরে সম্ভবত উৎসাহী কমিউনিস্ট
হতেন, নিজের দেশের বিশ্বাসঘাতক হতেন না। স্টালনেকার--লুইসের
ব্যাখ্যাও এই স্বজ্ঞাগত সত্যমানগুলি সমর্থন করে। আমাদের জগতের
তুলনায় হুভারের জন্মস্থান ও তার প্রাসঙ্গিক পরিণতি বদলানো,
অথচ অপ্রাসঙ্গিক অতিরিক্ত পরিবর্তন বাদ দেওয়া নিকটতম সম্ভাব্য জগতে তিনি
% BN-SRC-772: ন্যূনতম পরিবর্তনে জন্মস্থান বদলের প্রাসঙ্গিক পরিণতিও বদলায়; শুধু এক তথ্যের আক্ষরিক বদল নয়।
সোভিয়েত ইউনিয়নের সুনাগরিক হয়ে ওঠেন। প্রথম পূর্বধারণা ও
সিদ্ধান্তের পূর্বপক্ষ সত্য এমন নিকটতম জগৎ এটি; সেখানে তিনি
কমিউনিস্ট পার্টির অনুগত সদস্য, তাই বিশ্বাসঘাতক নন। কিন্তু
দ্বিতীয় পূর্বধারণা মূল্যায়নের জন্য অন্য জগৎ দেখতে হয়: হুভার
কমিউনিস্ট এমন নিকটতম জগৎটি। সেখানে তিনি যুক্তরাষ্ট্রেই জন্মান,
পরে মত বদলে কমিউনিস্ট হন এবং সেই অর্থে বিশ্বাসঘাতক হন।\footnote{উদাহরণটির
  তাৎপর্য বুঝতে কিছু অধিবিদ্যাগত ও রাজনৈতিক অনুমান মেনে নিতে হয়;
  যেমন, হুভার রুশ বাবা-মায়ের ঘরে জন্মাতে পারতেন, অথবা
  ১৯৫০-এর দশকে যুক্তরাষ্ট্রের কমিউনিস্টরা নিজেদের দেশের
  বিশ্বাসঘাতক ছিলেন।}

\begin{prob}
  প্রতিবাস্তব শর্তবচনের পরিযায়িতা ব্যর্থ হওয়ার একটি
  বিশ্বাসযোগ্য ও স্বজ্ঞাগত উদাহরণ দিন।
\end{prob}

\begin{ex}\ollabel{ex:trans-counterex}
  গোলক অর্থতত্ত্বে অনুমিতিটি অবৈধ, অর্থাৎ $p
  \cif q, q \cif r \Entails/ p \cif r$। ধরা যাক,
  $\mModel{M} = \tuple{W, O, V}$ মডেলে $W = \{w, w_1, w_2\}$,
  $O_w = \{\{w\}, \{w, w_1\}, \{w, w_1, w_2\}\}$,
  $V(p) = \{w_2\}$, $V(q) = \{w_1, w_2\}$
  এবং $V(r) = \{w_1\}$।
  % BN-SRC-774: অন্য দুই জগতের singleton গোলক দিয়ে O মোট করা যায়।
  অন্য দুই জগতে প্রতিটির নিজের এককবিশিষ্ট সেটমাত্র দিয়ে
  গোলকব্যবস্থা নেওয়া যায়। $p$-স্বীকারক গোলক
  $S = \{w, w_1, w_2\}$-এর প্রতিটি জগতে $p \lif q$
  সত্য; তাই $\mSat{M}{p \cif q}[w]$। আবার $q$-স্বীকারক গোলক
  % BN-SRC-426: এই গোলকের সর্বত্র q→r সত্য; নঞর্থক পরিতৃপ্তি নয়।
  % BN-SRC-775: একটি গোলকের প্রতিটি জগতে সূত্র সত্য; মডেলব্যাপী mSat নয়।
  $S' = \{w, w_1\}$-এর প্রতিটি জগতে $q \lif r$
  সত্য; তাই $\mSat{M}{q \cif r}[w]$। কিন্তু $p$-স্বীকারক
  গোলক $\{w, w_1, w_2\}$-এ এমন একটি জগৎ আছে—$w_2$—যেখানে
  $\mSat/{M}{p \lif r}[w_2]$। ওই পূর্বপক্ষ সত্য
  কেবল এই জগতে; তাই একমাত্র পূর্বপক্ষ-স্বীকারক গোলকেই
  শর্তমূলক অসত্য, এবং সিদ্ধান্তটি মিথ্যা।
\end{ex}

\begin{prob}
  \olref[cnt][min][tra]{ex:trans-counterex}-এর প্রত্যুদাহরণের
  সংশ্লিষ্ট গোলকচিত্র আঁকুন।
\end{prob}

\begin{prob}
  \olref[cnt][min][tra]{ex:trans-counterex}-এ জগৎ $w_2$-তে
  হুভার রাশিয়ায় জন্মান, কমিউনিস্ট হন, কিন্তু বিশ্বাসঘাতক নন;
  জগৎ $w_1$-তে তিনি যুক্তরাষ্ট্রে জন্মান, কমিউনিস্ট হন এবং
  বিশ্বাসঘাতক হন। এই মডেলে $w_1$ জগৎটি~$w$-র কাছে
  $w_2$-এর তুলনায় বেশি নিকটবর্তী। এটি কি অপরিহার্য?
  হুভারের রাশিয়ায় জন্মানোর সম্ভাবনাকে তাঁর কমিউনিস্ট
  হওয়ার সম্ভাবনার চেয়ে দূরবর্তী না ধরেও কি প্রত্যুদাহরণ
  দেওয়া যায়?
\end{prob}

\end{document}
